% 正文区：参赛项目说明书主体内容
% 篇幅建议 3-5 页
% ============================================

% ========== 元数据表格（与原始格式一致）==========
\begin{longtable}{|p{3cm}|p{10.5cm}|}
    \hline
    参赛题目类型（题目5、6、9填写） &
        \ifx\projectSubType\empty
            \textcolor{gray}{（如非题目5/6/9，可留空）}
        \else
            \projectSubType{}
        \fi \\
    \hline
    题目名称 & \projectTitle{} \\
    \hline
    \multicolumn{2}{|p{13.8cm}|}{%
        （以下内容可加页，篇幅3--5页）

        \textbf{设计原理与方法：}

        \textbf{实验仪器与装置：}

        \textbf{数据测量与分析：}

        \textbf{结论：}

        \textbf{项目创新与特色：}

        \textbf{制作成本（明细）：}
    } \\
    \hline
\end{longtable}

\vspace{0.5cm}

% ========== 以下为各板块详细撰写区 ==========
% 可根据需要取消注释，逐板块撰写；也可直接在表格中填入内容

% ---------- 设计原理与方法 ----------
\section{设计原理与方法}

在此处详细说明项目的物理原理、设计思路和实验方法。包括理论推导、关键公式、设计方案的选择依据等。

% ---------- 实验仪器与装置 ----------
\section{实验仪器与装置}

在此处列出实验所用仪器设备的名称、型号、主要参数，以及自行搭建的实验装置。可配图说明。

% ---------- 数据测量与分析 ----------
\section{数据测量与分析}

在此处展示实验测量数据，进行误差分析、数据处理和结果讨论。可使用表格和图表辅助说明。

% ---------- 结论 ----------
\section{结论}

在此处总结实验的主要发现和结论，评价实验结果与理论预期的符合程度。

% ---------- 项目创新与特色 ----------
\section{项目创新与特色}

在此处说明本项目的创新之处和特色亮点，与现有方法/方案的对比优势。

% ---------- 制作成本（明细）----------
\section{制作成本（明细）}

\begin{tabular}{|c|c|c|c|}
    \hline
    序号 & 名称 & 数量 & 单价（元） \\
    \hline
    1 &  &  &  \\
    \hline
    2 &  &  &  \\
    \hline
    3 &  &  &  \\
    \hline
    4 &  &  &  \\
    \hline
    5 &  &  &  \\
    \hline
    \multicolumn{3}{|r|}{合计} &  \\
    \hline
\end{tabular}

\vspace{1cm}

% ========== 注意事项（与原始格式一致）==========
{\heiti 注意：}

竞赛采用匿名评审，说明书里不允许出现任何参赛队和参赛学校信息。

\vspace{0.5cm}

% ========== 联系方式 ==========
\begin{center}
    {\songti\zihao{-5}
    邮箱: wlsyzx@bjtu.edu.cn \quad 联系电话: 51683629，13811861683 贺敬文

    通信地址：北京交通大学物理国家级实验教学示范中心 \quad 邮编: 100044
    }
\end{center}
